\documentclass[10pt,a4paper]{amsart}
\usepackage[margin=19mm]{geometry}
\usepackage{amsmath,amssymb,mathtools}
\usepackage[scr=rsfs,cal=euler]{mathalfa}
\usepackage{xfrac}
\usepackage[unicode,hidelinks,bookmarks=false]{hyperref}
\newcommand{\ie}{i.e.\ }
\newcommand{\cf}{cf.\ }
\newcommand{\resp}{respectively\ }
\newcommand{\Subsection}{\subsection}
\providecommand{\nobreakdash}{\nobreak\hbox{-}}
\setlength{\emergencystretch}{2em}
\multlinegap0pt
\usepackage{bm}
\usepackage{bbm}
\usepackage{dsfont}
\usepackage{xspace}
\usepackage{cancel}
\usepackage{mathtools}
\usepackage{amsthm}
\makeatletter
\@ifundefined{theorem}{\newtheorem{theorem}{Theorem}}{}
\@ifundefined{prop}{\newtheorem{prop}[theorem]{Proposition}}{}
\makeatother
\title[Generalised spin Calogero-Moser systems from Cherednik algeb]{Generalised spin Calogero-Moser systems from Cherednik algebras}
\author{Misha Feigin, Mikhail Vasilev, Martin Vrabec}
\date{Source: arXiv:2509.14099v2 (CC BY 4.0)}
\begin{document}
\maketitle

\noindent\emph{Source and licence.} Generalised spin Calogero-Moser systems from Cherednik algebras. Version arXiv:2509.14099v2 (CC BY 4.0), distributed under the Creative Commons Attribution 4.0 International licence, as recorded in the arXiv licence metadata for this exact version and shown on the abstract page https://arxiv.org/abs/2509.14099v2. Full rights notice: licenses/arxiv-2509.14099-CC-BY-4.0.txt.\par\smallskip

\noindent\textit{Editorial scope.} Counted unit: the Prop whose source label is \texttt{projref} in the article (source0.tex lines 386--396, source label \texttt{projref}), delivered together with the article's own complete original proof of it (source0.tex lines 398--435, one \texttt{proof} environment of 38 lines). No mathematics is written, repaired or re-derived: statement and proof are the article's own text line for line. Required context retained uncounted: none. Every definition, display and label of those lines is retained unchanged. Omitted from this package and not redistributed: 1--385, 397--397, 436--2366. Nothing inside a retained excerpt is omitted or altered: every retained body is delivered exactly as the article's lines are written, so no omission receipt of any kind accompanies this package. Citations: the retained excerpts carry no citation command at all, so no bibliography entry of the article is transcribed and no reference is redistributed. Reference adaptation: none was needed and none was made; every reference inside the delivered unit resolves inside it, so no unresolved reference or citation is produced. Printed slips kept literal: no printed word, symbol, hypothesis or number of the counted statement or of its proof was altered. Layout and class: the article's own class is \texttt{amsart} with packages including graphicx, amsmath, amssymb, mathtools, mathrsfs, tikz-cd, float, tikz, color, hyperref, longtable; the wrapper is the lane's pinned \texttt{amsart} editorial wrapper at 10pt on A4, which restates the theorem-like environments the retained text needs and adds the article's own macro definitions that the retained text needs (none), with extra wrapper packages (bm, bbm, dsfont, xspace, cancel, mathtools). The delivered numbering is the unit's own. Assets: no retained span contains a figure environment, a graphics-inclusion command, a PDF-page inclusion, a table or a TikZ picture; no image file is copied into this package, the package declares no asset dependency and no image file is redistributed with it. Licence: version arXiv:2509.14099v2 of this article is distributed under the Creative Commons Attribution 4.0 International licence, as recorded in the arXiv licence metadata for this exact version and shown on the abstract page https://arxiv.org/abs/2509.14099v2. A full rights notice is delivered at licenses/arxiv-2509.14099-CC-BY-4.0.txt. Characters: every retained excerpt is pure ASCII, so the delivered engine, XeLaTeX with its default Latin Modern text font, reports no missing character.
\begin{prop} \label{projref}
Let $U = V$ be the reflection representation of the Coxeter group $W$. 
Suppose that $\widehat{c}_{\widehat{\alpha}} \neq 0$ for  $\widehat{\alpha} \in \widehat{R}_+ \setminus \{0\}$. Then the operator
\begin{equation}
    \widehat{P}_{\widehat{\alpha}} = 1 + \frac{1}{\widehat{c}_{\widehat{\alpha}} (\widehat{\alpha}, \widehat{\alpha})} \sum_{\substack{ \gamma \in R_+ \\\widehat{\gamma} = \widehat{\alpha}}} 
    c_{\gamma}(\gamma, \gamma) (P_{\gamma} - 1)  \label{P hat}
\end{equation}
preserves
the space $\pi$, and when acting on $\pi$, it is equal to the orthogonal reflection 
about the hyperplane orthogonal to $\widehat{\alpha}$.
\end{prop}

\begin{proof}
Let us define 
\begin{equation*}
    \widetilde{P}_{\widehat{\alpha}} = \sum_{\substack{ \gamma \in R_+ \\ \widehat{\gamma} = \widehat{\alpha}}} 
     c_{\gamma}(\gamma, \gamma) P_{\gamma}.
\end{equation*}
For any $y \in \pi$, we have
\begin{equation}
\label{refl}
    \widetilde{P}_{\widehat{\alpha}}(y) %= \sum\limits_{\substack{\gamma \in R_{+} \\\widehat{\gamma} = \widehat{\alpha}}}c_{\gamma}(\gamma, \gamma) P_{\gamma}(y) 
    = \sum\limits_{\substack{\gamma \in R_{+} \\ \widehat{\gamma} = \widehat{\alpha}}}c_{\gamma}(\gamma, \gamma) y - 2\sum\limits_{\substack{\gamma \in R_{+} \\\widehat{\gamma} = \widehat{\alpha}}} c_{\gamma} (\gamma,y) \gamma. 
\end{equation}
To prove that the vector \eqref{refl} belongs to the subspace $\pi$, we need to simplify the second sum in \eqref{refl}. We claim that 
\begin{equation}
\label{elem}
     \sum\limits_{\substack{\gamma \in R_{+} \\ \widehat{\gamma} = \widehat{\alpha}}} c_{\gamma} (\gamma,y) \gamma =  \widehat{c}_{\widehat{\alpha}} (\widehat{\alpha},y) \widehat{\alpha}.
\end{equation}
Indeed, 
the set $S_{y} = \{c_{\gamma} (\gamma,y)\gamma \mid \gamma \in R_+, \, \widehat{\gamma} = \widehat{\alpha}\}$  
is $W_0$-invariant, where we use that $W_0$ is generated by reflections about the hyperplanes orthogonal to simple roots~$\Gamma_0^v$. 
Thus, the sum in the left-hand side of \eqref{elem} is fixed by $W_0$, hence belongs to the subspace $\pi$, and relation \eqref{elem} follows.
% Therefore, equality~\eqref{refl} becomes
% \begin{equation}
% \label{result1}
%     \widetilde{P}_{\widehat{\alpha}}(y) = \sum\limits_{\substack{\gamma \in R_{+} \\ \widehat{\gamma} = \widehat{\alpha}}}c_{\gamma}(\gamma, \gamma)  y - 2 \widehat{c}_{\widehat{\alpha}} (\widehat{\alpha},y) \widehat{\alpha}.
% \end{equation}

Now,
by using formula~\eqref{P hat} and equalities~\eqref{refl} and \eqref{elem},
we get that the action of the operator $\widehat{P}_{\widehat{\alpha}}$ on $\pi$ is
\begin{equation*}
    \widehat{P}_{\widehat{\alpha}} (y) = y +
    \frac{1}{\widehat{c}_{\widehat{\alpha}} (\widehat{\alpha}, \widehat{\alpha})} \bigg(  \widetilde{P}_{\widehat{\alpha}}(y)-\sum_{\substack{ \gamma \in R_+ \\\widehat{\gamma} = \widehat{\alpha}}} 
    c_{\gamma}(\gamma, \gamma) y \bigg)
    = y - \frac{2 (\widehat{\alpha}, y)}{(\widehat{\alpha}, \widehat{\alpha})} \widehat{\alpha},
\end{equation*}
which is the formula for the reflection on $\pi$ about the hyperplane orthogonal to~$\widehat{\alpha}$.
\end{proof}

\end{document}
